% Original template by Sebastian Del Bano Rollin, last updated by Oscar Bandtlow February 2024
\documentclass[12pt, oneside]{book}
\usepackage{graphicx}  % this is for includegraphics

\usepackage{setspace}
\onehalfspacing % this sets spacing to 1.5 


\usepackage{amsmath}
\usepackage{amsthm} % for theorems, lemmas etc
\usepackage{amsfonts}
\usepackage{lipsum} % to generate the lipsum random text in the sample


\usepackage[colorlinks=true, urlcolor=blue, pdfborder={0 0 0}]{hyperref}
\hypersetup{
     colorlinks   = true,
     citecolor    = blue
}

\theoremstyle{plain}
\newtheorem{theorem}{Theorem}[section]
\newtheorem{proposition}[theorem]{Proposition}
\newtheorem{lemma}[theorem]{Lemma}
\newtheorem{corollary}[theorem]{Corollary}
\newtheorem{fact}[theorem]{Fact}

\theoremstyle{definition}
\newtheorem{definition}[theorem]{Definition}
\newtheorem{example}[theorem]{Example}
\newtheorem{remark}[theorem]{Remark}
\newtheorem{remarks}[theorem]{Remarks}

\newcommand{\Cov}{\mathrm{Cov}}
\newcommand{\Var}{\mathrm{Var}}

\begin{document}

\begin{titlepage}
\begin{center}
        \vspace{-2cm}
Data Analytics MSc Dissertation MTHM038, 2023/24
%Financial Mathematics MSc Dissertation MTHM038, 2022/23
%Financial Computing MSc Dissertation MTHM038, 2022/23
%Mathematics MSc Dissertation MMTHM038, 2022/23
		\\
        \Huge
        \textbf{Constructibility of polygons}
        \\        
        \vspace{0.4cm}
        \LARGE
        With special emphasis on examples
        \\        
        \vspace{0.4cm}        
        \textbf{Carl Friedrich Gauss, ID 036782976}% student name and number        
        \\
        \large Supervisor: Prof.\ Johann Friedrich Pfaff
        \\
        \vspace{0.9cm}
        \includegraphics[scale=0.3]{QMCrest.jpg}
        \\
        \vspace{0.9cm}        
        \LARGE 
        A thesis presented for the degree of\\
        Master of Science in \emph{name of programme}\\
        \vspace{0.7cm}        
        \Large
        School of Mathematical Sciences\\ 
        and \emph{possible joint school}\\
        Queen Mary University of London \\
    \end{center}
\end{titlepage}


\chapter*{Declaration of original work}
\begin{flushright}
This declaration is made on \today.
\end{flushright}


{\bf Student's Declaration:}
I Student Name hereby declare that the work in this thesis 
is my original work. I have not copied from any other students' work, 
work of mine submitted elsewhere,  or from any other sources except where due 
reference or acknowledgement is made explicitly in the text. 
Furthermore, no part of this dissertation has been written for me by another person,
by generative artificial intelligence (AI), or by AI-assisted technologies.  
 



Referenced text has been flagged by:
\begin{enumerate}
\item Using italic fonts, {\bf and} % LaTeX: {\it text}  
\item using quotation marks ``\ldots '', {\bf and}
\item explicitly mentioning the source in the text.
\end{enumerate}

%This excludes any definitions known from your modules or that can be found in an undergraduate text book.

\newpage

\thispagestyle{empty}
        \begin{flushright}
                This work is dedicated to my dog Charles Frederick.
        \end{flushright}
\vspace{\stretch{2}}\null



\chapter*{Acknowledgements}
\lipsum[100] % replace this by your text

\chapter*{Abstract}
\lipsum[100]% replace this by your text

\tableofcontents

\chapter{Introduction}
This note presents a conjecture stemming from our investigations in the generation of sigmoid tensor categories of Picard numbers of tori in Banach algebras.
\lipsum[20]
 
\section{Motivation for this work}
In the works of Petri (\cite[Theorem 2.3]{Petri}) we find the following statement

\begin{theorem}[{\cite[Theorem 2.3]{Petri}, see also \cite[pg. 45]{BlackScholes}}]
The Gramm matrix for $E_8$ is:
$$
\begin{pmatrix}
2	&-1&0	&0	&0	&0	&0	&0\\
-1	&2	&-1	&0	&0	&0	&0	&0\\
0	&-1	&2	&-1	&0	&0	&0	&-1\\
0	&0	&-1	&2	&-1	&0	&0	&0\\
0	&0	&0	&-1	&2	&-1	&0	&0\\
0	&0	&0	&0	&-1	&2	&-1	&0\\
0	&0	&0	&0	&0	&-1	&2	&0\\
0	&0	&-1	&0	&0	&0	&0	&2
\end{pmatrix}.
$$
\end{theorem}

\subsection{The problem of exponential extensions}
\lipsum[23]
\subsection{The approach of Junderstein}
\lipsum[23]

\chapter{Eulerian topological string motives}
\lipsum[23]
\section{Definitions}
\lipsum[23]
\subsection{Tate's theorem}
\paragraph{Preliminary considerations}
\lipsum[23]
\paragraph{Motivic financial algebroids}
\lipsum[23]

\subsection{Grothendieck topologies}
\lipsum[23]

\section{Calculation of the invariant cycles}
\lipsum[23]

\subsection{Fontaine's theorem}
\lipsum[23]

\chapter{Conclusions}
\lipsum[100]

\appendix
\chapter{Implementation of the {\tt BarrierOptionCVA} class}
\lipsum[100]
\chapter[Shorter running title]{Additional details on the Gundermanian determinant}
\lipsum[100]



\bibliographystyle{plain}
\begin{thebibliography}{OOO}
\bibitem[P99]{Petri}
   William Petri, 
  {\it Analysis of infinitely generated frog complexes},
	Rendicoti Ran\ae \ Analysorum, 234 {\bf (4)}, 34--21, 2015

\bibitem[Ross]{Ross}
  Sheldon Ross, {\it An Elementary Introduction to Mathematical Finance},
	3rd Edition, Cambridge University Press, 2011
	
\bibitem[Hull]{Hull}
	John C. Hull, 
	{\it Options, Futures, and Other Derivatives},
	8th Edition, Pearson Education, 2011

\bibitem[BS]{BlackScholes}
	Fischer Black and  Myron Scholes,
	{\it \href{https://www.cs.princeton.edu/courses/archive/fall09/cos323/papers/black_scholes73.pdf}
	{The Pricing of Options and Corporate Liabilities}},
	Journal of Political Economy 81 {\bf 3}, 637--654,  (1973)
	
\end{thebibliography}



\end{document}
